\documentclass[twoside,12pt]{book}

\usepackage[utf8]{inputenc}
\usepackage[T1]{fontenc}
\usepackage{mathpazo}
\usepackage{amsmath,amssymb,amsthm}
\usepackage{graphicx}
\usepackage{listings}
\usepackage{xcolor}
\usepackage{enumitem}
\usepackage{tikz}
\usepackage{algorithm}
\usepackage{algpseudocode}
\usepackage{booktabs}
\usepackage{geometry}
\usepackage{fancyhdr}
\usepackage{titlesec}
\usepackage[colorlinks=true, linkcolor=blue, citecolor=blue, urlcolor=blue, plainpages=false, pdfpagelabels]{hyperref}

\geometry{
  papersize={6in,9in},
  margin=0.75in,
  top=0.75in,bottom=0.75in,
  inner=0.85in,outer=0.65in,
  headheight=14pt
}

\pagestyle{fancy}
\fancyhf{}
\fancyhead[LE]{\small\textit{11. Graph Algorithms}}
\fancyhead[RO]{\small\textit{Randomized Algorithms}}
\fancyfoot[C]{\thepage}
\renewcommand{\headrulewidth}{0.4pt}

\definecolor{codegray}{rgb}{0.45,0.45,0.45}
\lstdefinestyle{cppstyle}{
  basicstyle=\ttfamily\scriptsize,
  keywordstyle=\bfseries,
  commentstyle=\itshape\color{codegray},
  numbers=left,
  numberstyle=\tiny\color{codegray},
  numbersep=8pt,
  frame=single,
  framerule=0.4pt,
  rulecolor=\color{codegray},
  backgroundcolor=\color{gray!5},
  breaklines=true,
  tabsize=4,
  showstringspaces=false,
  language=C++,
  morekeywords={nullptr,size_t,auto,const,constexpr,
    unordered_map,vector,pair,mt19937,
    uniform_int_distribution,INT_MAX}
}
\lstset{style=cppstyle}

\theoremstyle{plain}
\newtheorem{theorem}{Theorem}[chapter]
\newtheorem{lemma}[theorem]{Lemma}
\newtheorem{corollary}[theorem]{Corollary}
\newtheorem{proposition}[theorem]{Proposition}
\theoremstyle{definition}
\newtheorem{definition}[theorem]{Definition}
\newtheorem{example}[theorem]{Example}
\newtheorem{remark}[theorem]{Remark}
\theoremstyle{remark}
\newtheorem{exercise}[theorem]{Exercise}
\newtheorem{problem}[theorem]{Problem}
\newenvironment{cpplisting}{\begin{center}\small}{\end{center}}

\newcommand{\Exp}{\operatorname{\mathbb{E}}}
\newcommand{\Var}{\operatorname{Var}}
\newcommand{\Cov}{\operatorname{Cov}}

\begin{document}

\chapter*{11 Graph Algorithms}
\addcontentsline{toc}{chapter}{11 Graph Algorithms}
\markboth{11 Graph Algorithms}{Randomized Algorithms}
\label{ch:graph-algorithms}

In this chapter we study several fundamental optimization problems on graphs: all-pairs shortest paths, minimum cuts, and minimum spanning trees. In each case, deterministic polynomial-time algorithms are known, but the use of randomization yields significantly faster solutions. The randomized algorithms we present exploit random sampling and random contraction to achieve running times that are difficult or impossible to obtain deterministically.

We begin with the all-pairs shortest paths problem, relating it to min-plus
and Boolean matrix multiplication and carefully separating safe baselines from
specialized faster reductions. We then study the minimum cut problem,
presenting Karger's elegant contraction algorithm and its dramatic improvement
due to Karger and Stein. Finally, we present the Karger--Klein--Tarjan
randomized linear-time algorithm for minimum spanning trees, which combines
random sampling with Bor\v{u}vka's algorithm and a linear-time verification
procedure.

%% ============================================================
\section{All-pairs Shortest Paths}
\label{sec:apsp}
%% ============================================================

The \emph{all-pairs shortest paths} (APSP) problem asks: given a weighted graph $G = (V, E)$ with $n$ vertices, compute the shortest-path distance $d(u, v)$ for every pair of vertices $u, v \in V$.

\begin{definition}[All-pairs shortest paths]
Given a graph $G = (V, E)$ with edge weights $w: E \to \mathbb{R}$, the APSP problem requires computing the $n \times n$ distance matrix $D$ where $D[u, v]$ is the length of the shortest path from $u$ to $v$. If $v$ is unreachable from $u$, then $D[u, v] = \infty$.
\end{definition}

The classical Floyd--Warshall algorithm solves APSP in $O(n^3)$ time. For sparse graphs, running Dijkstra from each vertex gives $O(nm + n^2 \log n)$ time. Can we do better?

\subsection{Min-plus matrix multiplication}

The key observation is the deep connection between shortest paths and a variant of matrix multiplication.

\begin{definition}[Min-plus product]
The \emph{min-plus product} (or \emph{distance product}) of two $n \times n$ matrices $A$ and $B$ is the $n \times n$ matrix $C$ defined by:
\[
C[i, j] = \min_{1 \leq k \leq n} \left( A[i, k] + B[k, j] \right).
\]
\end{definition}

Compare this with standard matrix multiplication, where we replace $+$ with $\times$ and $\min$ with $\sum$. The min-plus product captures the essence of shortest-path computation: the shortest path from $u$ to $v$ using at most $2^k$ edges can be expressed as a min-plus product of the distance matrices for paths of length at most $2^{k-1}$.

\begin{theorem}[APSP via repeated squaring]
\label{thm:apsp-squaring}
Let $W$ be the $n \times n$ weight matrix of a graph $G$ (with $W[i,j] = w(i,j)$ if $(i,j) \in E$ and $W[i,j] = \infty$ otherwise). Define the min-plus power $W^{(k)}$ recursively by:
\[
W^{(1)} = W, \qquad W^{(2k)} = W^{(k)} \star W^{(k)},
\]
where $\star$ denotes the min-plus product. Then $W^{(n-1)}$ is the all-pairs shortest path distance matrix.
\end{theorem}

\begin{proof}
By induction on $k$. The matrix $W^{(k)}[i,j]$ gives the shortest path from $i$ to $j$ using at most $k$ edges. For $k = 1$, this is just the direct edge weights. For the inductive step, a path of at most $2k$ edges can be split into two halves of at most $k$ edges each, so $W^{(2k)}[i,j] = \min_k (W^{(k)}[i,k] + W^{(k)}[k,j]) = (W^{(k)} \star W^{(k)})[i,j]$. Since the shortest path in an $n$-vertex graph uses at most $n - 1$ edges, $W^{(n-1)}$ gives the correct distances.
\end{proof}

This gives an $O(n^3 \log n)$ algorithm by computing $\lceil \log_2(n-1) \rceil$ min-plus products, each taking $O(n^3)$ time by brute force. But can we compute the min-plus product faster?

\subsection{Min-plus products and the scope of algebraic reductions}

The min-plus product resembles ordinary matrix multiplication, but the two
operations are not interchangeable. In particular, the tempting finite-field
encoding of entries as powers does not compute a min-plus product: a finite
field has no order with respect to which one summand can ``dominate'' the
others, and modular cancellation can make a shortest witness disappear.

\begin{theorem}[Safe min-plus baseline]
\label{thm:min-plus-baseline}
For arbitrary $n\times n$ matrices over the extended reals, assuming no
negative cycle is relevant to the requested distances, the min-plus product
can be computed in $O(n^3)$ time and $O(n^2)$ space. Repeated squaring
computes APSP in $O(n^3\log n)$ time.
\end{theorem}

\begin{proof}
For each pair $(i,j)$, scan all $k\in\{1,\ldots,n\}$ and retain the
minimum of $A[i,k]+B[k,j]$, treating $\infty$ in the usual way. There are
$n^2$ outputs and $n$ candidates per output, so the running time is
$O(n^3)$ and the output requires $O(n^2)$ space. The repeated-squaring
theorem above requires $\lceil\log_2(n-1)\rceil$ such products.
\end{proof}

Faster algorithms for restricted integer-weight instances use substantially
more structure, such as bounded weights, word-level operations, or carefully
analyzed algebraic encodings. They are not consequences of the invalid
finite-field dominance argument. The cited Alon--Galil--Margalit work is
therefore an advanced reference rather than a general black-box reduction in
this chapter. Whenever a faster min-plus routine is available under stated
additional assumptions, it may replace the baseline in the APSP framework.

\subsection{Seidel's algorithm for unweighted undirected graphs}

For unweighted, undirected graphs, Seidel's algorithm uses Boolean matrix
multiplication to reduce the distance scale by a factor of two at each
recursive level. It is deterministic once the matrix multiplication routine
is fixed; randomization is not required for the recurrence stated here.

\begin{theorem}[Seidel's APSP]
\label{thm:seidel-apsp}
Let $G$ be an unweighted undirected graph on $n$ vertices. The all-pairs shortest path distances can be computed in $O(M(n) \log n)$ expected time, where $M(n)$ is the time to multiply two $n \times n$ matrices.
\end{theorem}

\begin{proof}
Let $A$ be the adjacency matrix of $G$, with Boolean arithmetic, and define
$B=A\lor(A\cdot A)$. Thus $B[i,j]=1$ exactly when $i$ and $j$ are equal,
adjacent, or have a common neighbor. Recursing on the graph represented by
$B$ reduces the relevant distance scale by a factor of two.

For a finite-distance pair $i\ne j$, let $c(i,j)$ be the number of vertices
$k$ with $A[i,k]=A[k,j]=1$. If $c(i,j)>0$, the original distance is
$2D'[i,j]$; if $c(i,j)=0$, it is $2D'[i,j]-1$, where $D'$ is the distance
matrix returned by the recursive call. Adjacent pairs are handled directly,
and disconnected pairs retain distance $\infty$. These formulas follow by
comparing whether a shortest path can be decomposed into two-edge segments;
the absence of a common neighbor is precisely the odd-distance case.

Boolean matrix multiplication computes $A\cdot A$ and all values $c(i,j)$ in
$O(M(n))$ time, while reconstruction costs $O(n^2)$. There are
$O(\log n)$ distance scales, so the total running time is
$O(M(n)\log n)$ and the working space is $O(n^2)$ after matrix reuse.
\end{proof}

\begin{remark}
With a matrix multiplication routine of exponent $\omega<2.373$, Seidel's
algorithm runs in $O(n^{\omega}\log n)$ time, which is significantly faster
than the $O(n^3)$ Floyd--Warshall baseline. This statement applies to
unweighted undirected graphs. For general weighted graphs, the safe baseline
in this chapter remains min-plus repeated squaring; faster results require
additional assumptions and more specialized algorithms.
\end{remark}

\lstinputlisting[
  language=C++,
  caption={All-pairs shortest paths: Floyd--Warshall baseline and min-plus methods.},
  label={lst:apsp}]{src/chapter11/apsp.h}

%% ============================================================
\section{The Min-Cut Problem}
\label{sec:min-cut}
%% ============================================================

Recall from Chapter~1 that the \emph{minimum cut} problem asks: given a connected undirected graph $G = (V, E)$, find a partition $(S, \bar{S})$ of $V$ that minimizes the number of edges $|\delta(S)|$ crossing the partition.

\begin{definition}[Minimum cut]
The \emph{minimum cut} (or \emph{global min-cut}) of a graph $G = (V, E)$ is a partition $(S, V \setminus S)$ with $S \neq \emptyset$ and $S \neq V$ that minimizes $|\delta(S)| = |\{(u,v) \in E : u \in S, v \notin S\}|$.
\end{definition}

The minimum cut can be found in $O(nm)$ time via max-flow computations: fix a vertex $s$ and compute the minimum $s$-$t$ cut for every other vertex $t$, then take the minimum. Karger's algorithm achieves the same result with a dramatically different and much simpler approach.

\subsection{Karger's contraction algorithm}

The idea is beautifully simple: repeatedly pick a random edge and \emph{contract} it (merge its two endpoints into a single supernode), until only two supernodes remain. The edges between these two supernodes form a cut.

\begin{algorithm}
\caption{Karger's Contraction Algorithm}
\label{alg:karger}
\begin{algorithmic}[1]
\Function{MinCut}{$G = (V, E)$}
    \While{$|V| > 2$}
        \State Pick an edge $(u, v) \in E$ uniformly at random
        \State Contract $(u, v)$: merge $u$ and $v$ into a single supernode
        \State Remove all self-loops
    \EndWhile
    \State \Return the edges between the two remaining supernodes
\EndFunction
\end{algorithmic}
\end{algorithm}

\begin{definition}[Edge contraction]
Contracting an edge $(u, v)$ in a multigraph $G = (V, E)$ produces a new graph $G / (u,v)$:
\begin{enumerate}[nosep]
    \item Replace $u$ and $v$ by a single new vertex $w$.
    \item Replace each edge $(u, x)$ or $(v, x)$ by an edge $(w, x)$.
    \item Remove all self-loops (edges of the form $(w, w)$).
\end{enumerate}
\end{definition}

The crucial observation is:

\begin{lemma}[Contraction preserves cuts]
\label{lem:contraction-preserves}
If $(S, \bar{S})$ is a cut in $G$, then it is also a cut in $G / (u,v)$ whenever $(u,v) \notin \delta(S)$. In particular, no contraction can decrease the size of the minimum cut.
\end{lemma}

\begin{proof}
When we contract an edge $(u, v)$ that does not cross the cut $(S, \bar{S})$, both $u$ and $v$ belong to the same side. The new supernode inherits all the edges from $u$ and $v$, including the edges that crossed the cut. No new edges cross the cut, and no existing crossing edges are lost (except self-loops, which by assumption do not cross the cut). Therefore $|\delta(S)|$ is preserved.
\end{proof}

\begin{theorem}[Success probability of Karger's algorithm]
\label{thm:karger-prob}
Fix a minimum cut $C$ of $G$ with $|C| = k$. The probability that Karger's algorithm outputs $C$ is at least $\binom{n}{2}^{-1} = \frac{2}{n(n-1)}$.
\end{theorem}

\begin{proof}
Let $E_i$ be the event that no edge of $C$ is contracted at step $i$, for $i = 1, 2, \ldots, n-2$. We want to bound $\Pr[E_1 \cap E_2 \cap \cdots \cap E_{n-2}]$.

At step $i$, suppose $n - i + 1$ supernodes remain. Since $C$ is a minimum cut, every vertex has degree at least $k$ (otherwise, a vertex of degree $< k$ would give a smaller cut). Thus the number of edges is at least $k(n - i + 1) / 2$. The number of edges in $C$ is exactly $k$. Therefore:
\[
\Pr[E_i \mid E_1 \cap \cdots \cap E_{i-1}] \geq 1 - \frac{k}{k(n - i + 1)/2} = 1 - \frac{2}{n - i + 1} = \frac{n - i - 1}{n - i + 1}.
\]

By the chain rule for conditional probability:
\begin{align*}
\Pr\left[\bigcap_{i=1}^{n-2} E_i\right]
&= \prod_{i=1}^{n-2} \Pr[E_i \mid E_1 \cap \cdots \cap E_{i-1}] \\
&\geq \prod_{i=1}^{n-2} \frac{n - i - 1}{n - i + 1} \\
&= \frac{n-2}{n} \cdot \frac{n-3}{n-1} \cdot \frac{n-4}{n-2} \cdots \frac{1}{3} \\
&= \frac{2}{n(n-1)}.
\end{align*}

The telescoping product leaves only the first two terms of the numerator and the last two terms of the denominator.
\end{proof}

\begin{theorem}[Running time of Karger's algorithm]
A single run of Karger's algorithm can be implemented in $O(n^2)$ time using adjacency matrices, or $O(m)$ time using adjacency lists. Running it $O(n^2 \log n)$ times gives a minimum cut with high probability, for a total running time of $O(n^4 \log n)$ or $O(mn^2 \log n)$.
\end{theorem}

\begin{proof}
Each contraction merges two vertices and removes self-loops, taking $O(n)$ time with adjacency matrices. We perform $n - 2$ contractions, giving $O(n^2)$ per run. The probability of failure in one run is at most $1 - 2/n^2$. After $t$ independent runs, the failure probability is at most $(1 - 2/n^2)^t \leq e^{-2t/n^2}$. Setting $t = cn^2 \ln n$ gives failure probability at most $n^{-c}$.
\end{proof}

\begin{corollary}
Any graph on $n$ vertices has at most $\binom{n}{2}$ minimum cuts.
\end{corollary}

\begin{proof}
Each of the $\binom{n}{2}$ possible minimum cuts must be found with probability at least $\frac{2}{n(n-1)}$ by a single run. Since the runs are independent, if there were more than $\binom{n}{2}$ minimum cuts, some cut would have to be found with probability less than $\frac{2}{n(n-1)}$ in expectation, contradicting Theorem~\ref{thm:karger-prob}.
\end{proof}

\subsection{The Karger--Stein algorithm}

Karger's original algorithm requires $\Omega(n^2)$ independent repetitions for high probability, giving $\Omega(n^4)$ running time. Karger and Stein observed that the early contractions are much safer than the late ones: the probability of not destroying the min-cut decreases sharply as the graph shrinks. By re-running only the ``dangerous'' later stages, we can dramatically improve the running time.

\begin{algorithm}
\caption{Karger--Stein Recursive Contraction}
\label{alg:karger-stein}
\begin{algorithmic}[1]
\Function{MinCutKS}{$G = (V, E)$}
    \If{$|V| \leq 6$}
        \State \Return brute-force min-cut on $G$
    \EndIf
    \State Contract edges until $\lceil |V| / \sqrt{2} \rceil + 1$ vertices remain; let the result be $G'$
    \State $(S_1, \bar{S}_1) \gets \text{MinCutKS}(G')$ \Comment{first recursive call}
    \State $(S_2, \bar{S}_2) \gets \text{MinCutKS}(G')$ \Comment{second recursive call (independent)}
    \State \Return the cut among $(S_1, \bar{S}_1)$ and $(S_2, \bar{S}_2)$ with fewer crossing edges
\EndFunction
\end{algorithmic}
\end{algorithm}

The crucial parameter is the ``stopping threshold'' $\lceil n / \sqrt{2} \rceil + 1$. At this point, the probability that the min-cut has survived is approximately $1/2$:

\begin{lemma}
\label{lem:ks-half}
When contracting from $n$ vertices down to $l = \lceil n / \sqrt{2} \rceil + 1$, the probability that a fixed min-cut $C$ is preserved is at least $1/2$.
\end{lemma}

\begin{proof}
By the telescoping product analysis of Theorem~\ref{thm:karger-prob}, the survival probability is:
\[
\prod_{i=0}^{n-l-1} \left(1 - \frac{2}{n - i}\right) = \frac{l(l-1)}{n(n-1)}.
\]
With $l = \lceil n/\sqrt{2} \rceil + 1 \approx n/\sqrt{2}$:
\[
\frac{l(l-1)}{n(n-1)} \approx \frac{n^2/2}{n^2} = \frac{1}{2}.
\]
\end{proof}

\begin{theorem}[Running time of Karger--Stein]
\label{thm:karger-stein-time}
The Karger--Stein algorithm runs in $O(n^2 \log n)$ expected time per invocation.
\end{theorem}

\begin{proof}
The contraction from $n$ vertices to $\lceil n/\sqrt{2} \rceil + 1$ vertices takes $O(n^2)$ time. The two recursive calls are on graphs with $O(n/\sqrt{2})$ vertices. The recurrence is:
\[
T(n) = 2T\!\left(\frac{n}{\sqrt{2}}\right) + O(n^2).
\]
By the Master Theorem (with $a = 2$, $b = \sqrt{2}$, and $f(n) = O(n^2)$): since $\log_{\sqrt{2}} 2 = 2$, we are in the boundary case, giving $T(n) = O(n^2 \log n)$.
\end{proof}

\begin{theorem}[Success probability of Karger--Stein]
\label{thm:ks-prob}
A single invocation of Karger--Stein outputs a minimum cut with probability $\Omega(1 / \log n)$.
\end{theorem}

\begin{proof}
Let $P(n)$ be the success probability on a graph with $n$ vertices. By Lemma~\ref{lem:ks-half}, the contraction to $n/\sqrt{2}$ vertices succeeds with probability $\geq 1/2$. Conditional on success, each of the two recursive calls succeeds independently with probability $P(n/\sqrt{2})$. The overall success probability satisfies:
\[
P(n) \geq \frac{1}{2} \left(1 - (1 - P(n/\sqrt{2}))^2\right).
\]

Let $Q(n) = 1 - P(n)$. Then:
\[
Q(n) \leq 1 - \frac{1}{2}(1 - Q(n/\sqrt{2})^2) = \frac{1}{2} + \frac{1}{2}Q(n/\sqrt{2})^2.
\]

The solution to this recurrence is $Q(n) \leq 1 - \Theta(1/\log n)$, hence $P(n) = \Omega(1/\log n)$.
\end{proof}

\begin{corollary}
Running the Karger--Stein algorithm $O(\log^2 n)$ times independently and taking the best result gives a minimum cut with high probability, in total time $O(n^2 \log^3 n)$.
\end{corollary}

\begin{proof}
The probability that all $O(\log^2 n)$ runs fail is at most $(1 - \Omega(1/\log n))^{O(\log^2 n)} \leq e^{-\Omega(\log n)} = n^{-\Omega(1)}$.
\end{proof}

\begin{remark}
The Karger--Stein algorithm is one of the most elegant examples of the ``trade time for probability'' paradigm in randomized algorithms. By observing that failures are concentrated in the later stages of the algorithm, we can re-run only the expensive part, achieving a dramatic speedup.
\end{remark}

\lstinputlisting[
  language=C++,
  caption={Karger's contraction algorithm and the Karger--Stein recursive variant.},
  label={lst:mincut}]{src/chapter11/min_cut.h}

%% ============================================================
\section{Minimum Spanning Trees}
\label{sec:mst}
%% ============================================================

The \emph{minimum spanning tree} (MST) problem asks: given a connected, weighted graph $G = (V, E)$, find a spanning tree of minimum total weight.

\begin{definition}[Minimum spanning tree]
A \emph{minimum spanning tree} of a connected graph $G = (V, E)$ with edge weights $w: E \to \mathbb{R}$ is a spanning tree $T \subseteq E$ minimizing $\sum_{e \in T} w(e)$.
\end{definition}

Classical deterministic algorithms (Kruskal's, Prim's, Bor\v{u}vka's) run in $O(m \log n)$ or $O(m + n \log n)$ time. The question we address is: can we do better?

The answer is yes, using randomization. The Karger--Klein--Tarjan algorithm achieves $O(m)$ expected time---\emph{linear} in the size of the input---which is optimal.

\subsection{Preliminaries}

We assume all edge weights are distinct (ties can be broken arbitrarily). Under this assumption, the MST is unique. Two fundamental properties underlie MST algorithms:

\begin{lemma}[Cut property]
For any nonempty proper subset $S \subset V$, the lightest edge crossing the cut $(S, \bar{S})$ belongs to the MST.
\end{lemma}

\begin{lemma}[Cycle property]
For any cycle $C$ in $G$, the heaviest edge on $C$ does not belong to the MST.
\end{lemma}

\subsection{Bor\v{u}vka's algorithm}

Bor\v{u}vka's algorithm is a classical linear-time algorithm for MSTs in the parallel setting. Each \emph{Bor\v{u}vka step}:
\begin{enumerate}[nosep]
    \item For each connected component, find the lightest edge leaving that component (the \emph{minimum outgoing edge}).
    \item Add all these edges to the MST.
    \item Contract the newly added edges, reducing the number of components by at least a factor of 2.
\end{enumerate}

\begin{lemma}
After $O(\log n)$ Bor\v{u}vka steps, the graph is contracted to a single component, and the MST is found.
\end{lemma}

Each Bor\v{u}vka step takes $O(m)$ time, giving $O(m \log n)$ total. The key insight of Karger, Klein, and Tarjan is to combine Bor\v{u}vka steps with random sampling to discard edges that provably cannot be in the MST.

\subsection{The sampling lemma}

The core of the Karger--Klein--Tarjan algorithm is a random-sampling lemma that bounds the number of ``useful'' edges after sampling.

\begin{definition}[$F$-heavy and $F$-light]
Let $F$ be a forest in $G$. An edge $e = \{u, v\}$ is \emph{$F$-heavy} if $w(e) > w_F(u, v)$, where $w_F(u, v)$ is the maximum weight edge on the unique path from $u$ to $v$ in $F$ (or $\infty$ if $u$ and $v$ are in different components of $F$). Otherwise $e$ is \emph{$F$-light}.
\end{definition}

By the cycle property, no $F$-heavy edge can be in the MST of $G$.

\begin{lemma}[Sampling lemma]
\label{lem:sampling}
Let $H$ be a random subgraph of $G$ obtained by including each edge independently with probability $p$, and let $F$ be the MST of $H$. Then the expected number of $F$-light edges in $G$ is at most $n/p$.
\end{lemma}

\begin{proof}
We analyze the process of building $H$ and $F$ simultaneously using a variant of Kruskal's algorithm. Process all edges of $G$ in increasing weight order. For each edge $e$:
\begin{enumerate}[nosep]
    \item If $e$ connects two different components of the current forest $F$, flip a biased coin: include $e$ in $H$ with probability $p$.
    \item If included, add $e$ to $F$.
    \item If $e$ connects vertices in the same component of $F$, it is $F$-heavy regardless of inclusion in $H$.
\end{enumerate}

The forest $F$ has at most $n - 1$ edges. Consider only the ``coin flips'' for edges that are $F$-light at the time of processing. Each such flip is a coin with probability $p$ of heads. The number of flips until $n - 1$ heads have occurred follows a negative binomial distribution with parameters $(n-1, p)$, which has expectation $(n-1)/p < n/p$.
\end{proof}

\subsection{The algorithm}

\begin{algorithm}
\caption{Karger--Klein--Tarjan MST Algorithm}
\label{alg:kkt}
\begin{algorithmic}[1]
\Function{FindMST}{$G = (V, E)$}
    \State Sample each edge of $G$ with probability $p = 1/2$ to form $H$
    \State $F \gets \text{FindMST}(H)$ \Comment{recurse on smaller graph}
    \State Find all $F$-light edges of $G$ \Comment{$O(m)$ time via verification}
    \State $G' \gets$ subgraph of $G$ induced by $F$-light edges
    \State Run $O(\log n)$ Bor\v{u}vka steps on $G'$, producing contracted graph $G''$
    \State \Return MST of $G''$ \Comment{recurse on smaller graph}
\EndFunction
\end{algorithmic}
\end{algorithm}

\begin{theorem}[Running time of KKT]
\label{thm:kkt}
The Karger--Klein--Tarjan algorithm finds a minimum spanning tree in $O(m)$ expected time.
\end{theorem}

\begin{proof}
Let $T(n, m)$ be the expected running time. By the sampling lemma, the expected number of $F$-light edges is at most $n/p = 2n$. After $O(\log n)$ Bor\v{u}vka steps on these $O(n)$ edges, the graph contracts to $O(n / 2^{\log n}) = O(1)$ vertices (a constant number of components), so the recursion bottoms out after $O(\log n)$ levels.

At each level, the total work is:
\begin{itemize}[nosep]
    \item Sampling: $O(m)$ to build $H$.
    \item Recursive MST on $H$: $T(n, m/2)$ in expectation (half the edges).
    \item Finding $F$-light edges: $O(m)$ using linear-time MST verification.
    \item Bor\v{u}vka steps: $O(m)$ on the $O(n)$ light edges.
    \item Recursive MST on contracted graph: $T(O(n), O(n))$.
\end{itemize}

The recursion tree has depth $O(\log n)$ (since each level halves the number of edges), and at each level the total work across all nodes is $O(m)$. Therefore $T(n, m) = O(m)$.
\end{proof}

\begin{remark}
The KKT algorithm is remarkable for several reasons:
\begin{enumerate}[nosep]
    \item It achieves the theoretical optimum of $O(m)$ time, matching the time needed just to read the input.
    \item It uses the comparison-only model (no bit manipulation tricks), unlike the Fredman--Willard algorithm.
    \item It combines three fundamentally different techniques: random sampling (to discard edges), Bor\v{u}vka's algorithm (to reduce the graph), and linear-time MST verification (to identify heavy edges).
\end{enumerate}
\end{remark}

\lstinputlisting[
  language=C++,
  caption={Karger--Klein--Tarjan randomized linear-time MST algorithm.},
  label={lst:mst}]{src/chapter11/mst.h}

%% ============================================================
\addcontentsline{toc}{section}{Notes}
\section*{Notes}
\label{sec:notes-10}
%% ============================================================

The study of randomized graph algorithms has produced some of the most beautiful results in the theory of computation. We highlight the key contributions behind the algorithms in this chapter.

\textbf{All-pairs shortest paths.} The connection between shortest paths and matrix multiplication was observed by several authors, including Shimbel (1955), Bellman (1958), and Leyzorek et al.\ (1957). Specialized algebraic reductions for restricted min-plus instances were developed by Fredman (1975) and Alon, Galil, and Margalit (1992); their hypotheses and encoding details are substantially more delicate than a finite-field ``dominant term'' argument. Seidel's $O(M(n) \log n)$ algorithm for unweighted undirected graphs appeared in 1992 and remains the standard matrix-multiplication-based treatment of this special case.

\textbf{Minimum cuts.} Karger's contraction algorithm was introduced in his 1993 paper ``Global min-cuts in RNC, and other ramifications of a simple min-cut algorithm.'' The analysis via telescoping products is elementary yet powerful. The Karger--Stein improvement appeared in their 1996 paper ``A new approach to the minimum cut problem,'' achieving the $O(n^2 \log^3 n)$ bound. The connection between min-cuts and network reliability, as well as the cactus representation of all minimum cuts, are developed extensively in Karger's subsequent work.

\textbf{Minimum spanning trees.} The linear-time MST verification problem was solved by Dixon, Rauch, and Tarjan (1992) and independently by King (1995). Karger (1993) showed how to use random sampling to discard edges, obtaining an $O(n \log n + m)$ algorithm. The full $O(m)$ randomized algorithm combining sampling, Bor\v{u}vka steps, and verification was achieved by Karger, Klein, and Tarjan in their 1995 paper ``A randomized linear-time algorithm to find minimum spanning trees.''

\begin{itemize}[nosep]
    \item R.~Seidel, ``On the All-Pairs-Shortest-Path Problem in Unweighted Undirected Graphs,'' \textit{J.\ Comput.\ Syst.\ Sci.}, 51(2):291--298, 1995.
    \item N.~Alon, Z.~Galil, and O.~Margalit, ``On the Exponent of the All Pairs Shortest Path Problem,'' \textit{J.\ Comput.\ Syst.\ Sci.}, 54(2):255--262, 1997.
    \item D.~R.~Karger, ``Global Min-Cuts in RNC, and Other Ramifications of a Simple Min-Cut Algorithm,'' \textit{Proc.\ 4th ACM--SIAM Symposium on Discrete Algorithms}, 21--30, 1993.
    \item D.~R.~Karger and C.~Stein, ``A New Approach to the Minimum Cut Problem,'' \textit{J.\ ACM}, 43(4):601--640, 1996.
    \item D.~R.~Karger, P.~N.~Klein, and R.~E.~Tarjan, ``A Randomized Linear-Time Algorithm to Find Minimum Spanning Trees,'' \textit{J.\ ACM}, 42(2):519--532, 1995.
    \item B.~Dixon, M.~Rauch, and R.~E.~Tarjan, ``Verification and Sensitivity Analysis of Minimum Spanning Trees in Linear Time,'' \textit{SIAM J.\ Comput.}, 21(6):1184--1192, 1992.
    \item V.~King, ``A Simpler Minimum Spanning Tree Verification Algorithm,'' \textit{Algorithmica}, 18:263--270, 1997.
    \item M.~L.~Fredman, ``New Bounds on Efficient Data Structures for Computing Minimum Spanning Trees,'' \textit{Proc.\ 15th ACM STOC}, 277--284, 1975.
\end{itemize}

%% ============================================================
\addcontentsline{toc}{section}{Problems}
\section*{Problems}
\label{sec:problems-10}
%% ============================================================

\begin{problem}[10.1]
Prove that the min-plus product of two $n \times n$ matrices can be computed in $O(n^3)$ time by brute force. Then show that computing the APSP distances via repeated min-plus squaring (Theorem~\ref{thm:apsp-squaring}) requires $O(n^3 \log n)$ time.
\end{problem}

\begin{problem}[10.2]
Let $G$ be an unweighted undirected graph with adjacency matrix $A$. Show that $(A^k)[i, j]$ (computed using standard arithmetic) equals the number of walks of length exactly $k$ from $i$ to $j$. How does this relate to shortest-path distances?
\end{problem}

\begin{problem}[10.3]
Implement Karger's contraction algorithm and test it on random graphs $G(n, p)$ with $n = 100$ and various values of $p$. Estimate the probability that a single run finds the minimum cut, and compare with the theoretical lower bound $\frac{2}{n(n-1)}$.
\end{problem}

\begin{problem}[10.4]
Prove that any graph on $n$ vertices has at most $\binom{n}{2}$ distinct minimum cuts. \textit{Hint:} Each minimum cut must be ``found'' by Karger's algorithm with probability at least $\frac{2}{n(n-1)}$.
\end{problem}

\begin{problem}[10.5]
Verify that the Karger--Stein recurrence $T(n) = 2T(n/\sqrt{2}) + O(n^2)$ solves to $O(n^2 \log n)$ using the Master Theorem. What is the precise constant in the $O$-notation?
\end{problem}

\begin{problem}[10.6]
In the Karger--Stein algorithm, show that the survival probability of a fixed min-cut when contracting from $n$ vertices to $\lceil n/\sqrt{2} \rceil + 1$ vertices is exactly $\frac{l(l-1)}{n(n-1)}$ where $l = \lceil n/\sqrt{2} \rceil + 1$.
\end{problem}

\begin{problem}[10.7]
Prove the cut property: for any nonempty proper subset $S \subset V$, the lightest edge crossing the cut $(S, \bar{S})$ belongs to the MST.
\end{problem}

\begin{problem}[10.8]
Prove the cycle property: for any cycle $C$ in $G$, the heaviest edge on $C$ does not belong to the MST.
\end{problem}

\begin{problem}[10.9]
In the sampling lemma (Lemma~\ref{lem:sampling}), explain why the negative binomial distribution arises. What is the exact expected number of $F$-light edges when $p = 1/2$?
\end{problem}

\begin{problem}[10.10]
\textbf{(Research)} Design a randomized algorithm for the minimum $s$-$t$ cut problem that runs in $O(n^2)$ expected time. \textit{Hint:} adapt Karger's contraction algorithm by contracting only edges that do not separate $s$ and $t$.
\end{problem}

\end{document}
